\chapter{充分性与分布变换习题}

本组对应第一次作业第 7--12 题。第 1--6 题在本轮学习讨论中没有完整解答，不计入已整理范围。以下重述数学问题并给出解答；几乎必然等式均指对每个模型分布分别成立，所用统计量和表示函数不含未知参数。

\section{柯西位置尺度族的最小充分统计量}
设 $X_1,\ldots,X_n$ 独立同分布于 $C(\mu,\sigma)$，$\mu\in\R$、$\sigma>0$ 均未知。证明全部次序统计量组成的向量最小充分。

\textbf{解。} 联合密度处处为正，且
\[
 p_{\mu,\sigma}(x)=\frac{\sigma^n}{\pi^n\prod_i[(x_i-\mu)^2+\sigma^2]}.
\]
排序不改变乘积，所以因子分解给出 $T=(X_{(1)},\ldots,X_{(n)})$ 充分。若两条似然成比例，令 $\sigma\to\infty$，两边分母同为首项系数 $1$ 的 $\sigma^{2n}$ 多项式，比例常数必为 $1$。于是对所有 $\mu\in\R$、$\sigma>0$，
\[
 \prod_i[(x_i-\mu)^2+\sigma^2]
 =\prod_i[(y_i-\mu)^2+\sigma^2].
\]
固定 $\mu$ 并令 $\sigma\downarrow0$，得 $P_x(\mu)^2=P_y(\mu)^2$，其中
$P_x(z)=\prod_i(z-x_i)$、$P_y(z)=\prod_i(z-y_i)$。因此多项式
$(P_x-P_y)(P_x+P_y)$ 恒为零。实多项式环无零因子，且两个多项式均首一，所以 $P_x=P_y$。实根及其重数相同，故排序后的样本相同。反之，排序相同显然使似然相同。由第二章的似然成比例判别，$T$ 最小充分。

\section{均匀位置族：最小充分但不完备}
设 $n\ge2$、$X_i\sim U(\theta,\theta+1)$ 独立同分布。证明 $T=(X_{(1)},X_{(n)})$ 最小充分而不完备。

\textbf{解。} 对正似然样本，允许的参数集合为
\[
 (x_{(n)}-1,x_{(1)}),\qquad
 p_\theta(x)=\ind_{\{x_{(n)}-1<\theta<x_{(1)}\}}.
\]
因子分解给出充分性。两条非零似然成比例，等价于上述非空开区间相同，因而两个端点分别相同，得到最小充分性。比较的是完整似然，不是可能出现 $0/0$ 的密度比。

令 $R=X_{(n)}-X_{(1)}$。由 $X_i=\theta+U_i$，$R$ 的分布不含 $\theta$。其密度及期望为
\[
 f_R(r)=n(n-1)r^{n-2}(1-r)\ind_{(0,1)}(r),\qquad
 \E R=\frac{n-1}{n+1}.
\]
于是固定有界函数 $g(T)=R-(n-1)/(n+1)$ 对所有参数均值为零，却不几乎必然为零。因此 $T$ 不有界完备，也不完备。第二章给出了该密度及 Basu 反证的详细推导。

\section{可以恢复充分统计量的统计量仍充分}
在本讲义的受支配模型框架内，设 $T$ 充分且 $T=\psi(S)$，$\psi$ 为不含参数的可测函数，证明 $S$ 充分。

\textbf{解。} 在受支配模型中，由因子分解，$p_\theta(x)=a_\theta(T(x))h(x)$。代入得
\[
 p_\theta(x)=a_\theta(\psi(S(x)))h(x).
\]
右端仍是 $S(x)$ 的函数与参数无关因子的乘积，故 $S$ 充分。方向不可颠倒：充分统计量的任意函数未必充分，例如把它映射成常数。

\section{完备充分统计量与最小充分统计量的关系}
设 $T$ 完备（或有界完备）且充分，最小充分统计量 $S$ 存在。证明 $T$ 最小充分，且 $S$ 也完备（或有界完备）。

\textbf{解。} 最小充分性先给出 $S=\psi(T)$。若 $a(S)$ 可积（在有界版本中要求有界）且所有参数下均值为零，则 $a(\psi(T))$ 具有相同性质；由 $T$ 的完备性，它几乎必然为零，所以 $S$ 完备。

为证明反向可恢复性，任取有界可测 $b$。$S$ 充分，故可统一选取不含参数的 $k$，使
\[
 k(S)=\E_\theta[b(T)\mid S].
\]
若 $|b|\le M$，可取 $|k|\le M$。残差 $b(T)-k(\psi(T))$ 是 $T$ 的固定有界函数，所有参数下均值为零，故由有界完备性
\[
 b(T)=k(S)\quad P_\theta\text{-a.s.}
\]
实值时取 $b(t)=\arctan t$，反解得到 $T=\tan(k(S))$；逆函数在有效像外可补定义。有限维统计量对每个坐标分别如此处理。一般标准 Borel 值域可用有界 Borel 编码。任意充分统计量 $U$ 都能恢复 $S$，从而恢复 $T$，所以 $T$ 最小充分。不能直接对未保证可积的 $T$ 作条件期望，也不能把一般无界残差用于有界完备性。

\section{平移指数分布的间距与独立性}
设 $n>2$，样本密度为
\[
 f_{a,\theta}(x)=\theta^{-1}e^{-(x-a)/\theta}\ind_{\{x>a\}},
 \quad a\in\R,\quad\theta>0.
\]
证明 $X_{(1)}$ 与 $Q=\sum_i(X_i-X_{(1)})$ 独立；再证明向量
\[
 Z_i=\frac{X_{(n)}-X_{(i)}}{X_{(n)}-X_{(n-1)}},\quad 1\le i\le n-2,
\]
与 $(X_{(1)},Q)$ 独立。

\textbf{解。} 定义缩放间距
\[
 W_1=n(X_{(1)}-a),\quad
 W_j=(n-j+1)(X_{(j)}-X_{(j-1)}),\quad 2\le j\le n.
\]
反变换为 $X_{(j)}=a+W_1/n+\sum_{k=2}^jW_k/(n-k+1)$，其 Jacobian 绝对值为 $1/n!$。次序统计量密度为
\[
 n!\theta^{-n}\exp\{-\textstyle\sum_j(x_j-a)/\theta\}
 \ind_{\{a<x_1<\cdots<x_n\}}.
\]
又有 $\sum_j(x_j-a)=\sum_jw_j$，故
\[
 f_W(w)=\prod_{j=1}^n\theta^{-1}e^{-w_j/\theta}\ind_{\{w_j>0\}}.
\]
即 $W_j$ 独立同分布于尺度为 $\theta$ 的指数分布。计数每个间距出现的次数，得
\[
 Q=\sum_{j=2}^nW_j,\qquad X_{(1)}=a+W_1/n,
\]
所以第一问成立。

令 $P_j=W_j/Q$，$j=2,\ldots,n$；只需 $P_2,\ldots,P_{n-1}$ 为自由变量。变换 $w_j=qp_j$ 的 Jacobian 为 $q^{n-2}$，于是
\[
 f_{Q,P}(q,p)=
 \frac{q^{n-2}e^{-q/\theta}}{(n-2)!\theta^{n-1}}\ind_{\{q>0\}}
 \cdot (n-2)!\ind_{\Delta}(p),
\]
其中 $\Delta=\{p_j>0,\sum_{j=2}^{n-1}p_j<1\}$。前者是 Gamma 密度，后者在体积为 $1/(n-2)!$ 的单纯形上积分为 $1$，故 $Q\perp\!\!\!\perp P$。再由 $W_1$ 与 $(W_2,\ldots,W_n)$ 独立，联合密度分解为
$f_{W_1}f_Qf_P$，得到 $P\perp\!\!\!\perp(W_1,Q)$。这是联合独立性，不能仅由两次两两独立直接推断。

最后
\[
 Z_i=\frac{\sum_{k=i+1}^nP_k/(n-k+1)}{P_n},\quad P_n>0\text{ a.s.}
\]
整个 $Z$ 向量是 $P$ 的函数，$(X_{(1)},Q)$ 是 $(W_1,Q)$ 的函数，故第二问成立。该结论不声称 $Z_i$ 之间相互独立。

\section{两段支撑模型的最小充分统计量}
设 $\theta>0$，$X_i$ 独立同分布，密度为
\[
 f_\theta(x)=\frac1{2\theta}
 [\ind_{(0,\theta)}(x)+\ind_{(2\theta,3\theta)}(x)].
\]
求一个最小充分统计量。

\textbf{解。} 只需考虑共同满测集 $(0,\infty)^n$。令
\[
\begin{aligned}
 M&=\max_i x_i,\\
 A&=\max(\{x_i:x_i<M/2\}\cup\{0\}),\\
 B&=\min\{x_i:x_i\ge M/2\},\\
 L&=\max(M/3,A),\qquad U=B/2.
\end{aligned}
\]
$B$ 总有定义，因为最大值属于相应集合。令 $(L^*,U^*)=(L,U)$ 若 $L<U$，否则令其为 $(0,0)$。下面证明 $T=(M,L^*,U^*)$ 最小充分。

由最大值所在区间，正似然参数必满足 $\theta>M$ 或 $M/3<\theta<M/2$。前者使全部观测落入第一段。后者中，$x_i<M/2$ 时只能属于第一段，要求 $x_i<\theta$；$x_i\ge M/2$ 时只能属于第二段，要求 $2\theta<x_i$，而 $x_i<3\theta$ 已由 $x_i\le M<3\theta$ 保证。故正似然参数集合恰为
\[
 C(x)=(M,\infty)\cup(L,U).
\]
空区间用 $(L^*,U^*)=(0,0)$ 统一编码。联合密度为
\[
 p_\theta(x)=(2\theta)^{-n}
 [\ind_{\{\theta>M\}}+\ind_{\{L^*<\theta<U^*\}}],
\]
给出充分性。

若 $p_\theta(x)=c\,p_\theta(y)$ 对所有参数成立，取 $\theta>\max(M_x,M_y)$，有 $c=1$，所以 $C(x)=C(y)$。该集合唯一无界连通分支的左端点恢复 $M$；剩余非空有界分支恢复 $L,U$，剩余为空则恢复编码 $(0,0)$。因此 $T(x)=T(y)$。反之 $T$ 相同使似然相同。由似然成比例判别，$T$ 最小充分。

\begin{notebox}[支撑与空区间]
记录区间端点时必须把所有空区间统一编码。否则不同的 $L\ge U$ 可能对应同一条似然，额外保留无用端点便不能由上述判别得到最小性。
\end{notebox}
